\documentclass[10pt]{article}
\usepackage[margin=0.95in]{geometry}
\usepackage{amsmath,amssymb,amsthm,microtype,xcolor}
\usepackage[colorlinks=true,urlcolor=blue!50!black]{hyperref}
\setlength{\emergencystretch}{3em}
\newtheorem{theorem}{Theorem}
\title{Reciprocal Summability of Every Unbounded\\Natural Collatz Orbit: A Lean Proof}
\author{Collatz proof project}
\date{September 5, 2026}
\begin{document}
\maketitle
\begin{abstract}
Every unbounded positive natural orbit of the shortcut Collatz map has a
convergent series of reciprocal values. We prove this in Lean by combining
an equal-time injectivity argument with existing binomial and growth bounds.
An explicit finite packing inequality supplies a geometric shell majorant.
Consequently every such orbit has bounded ideal correction, and its inverse
multiplicative drift tends to zero. This closes the correction-coverage gap
previously recorded in this repository; it does not exclude unbounded
orbits, prove a strict density gap, or rule out nontrivial cycles. We make
no claim of mathematical novelty.
\end{abstract}

\section{Statements and conventions}
Let $T(n)=n/2$ for even $n$ and $(3n+1)/2$ for odd $n$.
Set $x_t=T^t(N)$, and let $j_t$ count odd values among $x_0,\ldots,x_{t-1}$.
Define the nonnegative ideal correction $c_t(N)$ by the exact identity
\[
x_t=\frac{3^{j_t}}{2^t}\bigl(N+c_t(N)\bigr).
\]
In the repository, $\operatorname{Supercritical}(N)$ means only that
$c_t(N)$ is bounded above as $t$ varies. It does not mean a strict density
gap above $\log 2/\log 3$.
\begin{theorem}
If $N>0$ and $\forall B\,\exists t\;(B<x_t)$, then
\[
\sum_{t\ge0}\frac1{x_t}<\infty,
\qquad \sup_t c_t(N)<\infty,
\qquad \frac{2^t}{3^{j_t}}\longrightarrow0.
\]
For positive seeds, bounded ideal correction is equivalent to unboundedness
of the natural orbit.
\end{theorem}
An unbounded deterministic orbit cannot repeat a state. Moreover,
$T^k$ is injective on its entire orbit set for every fixed $k$: coalescence
at equal times would imply a repetition at two later times.

\section{Finite packing with explicit constants}
\begin{theorem}
Let $S\subseteq[0,32^m)\cap\mathbb N$ be finite, and suppose $T^{5m}$ is
injective on $S$. Then
\[
8^m|S|\le2\cdot216^m+243^m.
\]
\end{theorem}
Put $k=5m$. For any $n<2^k$, the existing upper master inequality gives
\[
2^k\bigl(T^k(n)+1\bigr)
 \le3^{j_k(n)}\bigl(n+2^{k-j_k(n)}\bigr),
\qquad T^k(n)<2\cdot3^{j_k(n)}.
\]
Split $S$ at odd count $3m$. Low-count elements map into
$[0,2\cdot27^m)$; injectivity bounds their number by $2\cdot27^m$.
The high-count elements are bounded by all residues with at least $3m$ odd
steps. The verified parity bijection gives exactly $\binom{k}{j}$ residues
with $j$ odd steps. Therefore the weighted binomial identity gives
\[
2^{3m}\sum_{j\ge3m}\binom{5m}{j}
\le\sum_j2^j\binom{5m}{j}=3^{5m}.
\]
Multiplying the low-count bound by $8^m$ and adding proves the inequality.
This applies to every finite subset of an unbounded orbit below $32^m$,
including values reached at arbitrarily late times.

\section{Geometric shells and time summability}
For orbit values in $[32^m,32^{m+1})$, the packing bound at scale $m+1$
and $1/n\le32^{-m}$ give
\[
\sum_{n\text{ in the shell}}\frac1n
\le b_m:=54\left(\frac{27}{32}\right)^m
 +\frac{243}{8}\left(\frac{243}{256}\right)^m.
\]
Both ratios are less than one. Inductively splitting any finite set of
positive orbit values at the shell endpoints bounds its reciprocal sum by
an initial partial sum of $\sum_m b_m$.

For a finite time prefix, take the image of its orbit values as a finite
set. Unboundedness supplies injectivity in time, so passing to this set
preserves the reciprocal sum without losing multiplicities. Each such sum
is bounded by $\sum_m b_m$, proving summability. The displayed majorant
has sum less than 1000; the quantitative refinement below bounds every
value tail explicitly.

This distinction between values and time is essential. A finite cycle has
a finite reciprocal sum over its distinct values but a divergent reciprocal
series over its infinite time sequence. The packing step alone would not
justify summability for a repeating orbit.

\section{Correction coverage and drift}
The earlier module \texttt{Reciprocal.lean} proves, for positive seeds,
\[
\operatorname{Supercritical}(N)
\quad\Longleftrightarrow\quad
\text{the time-indexed reciprocal series is summable}.
\]
It also proves that bounded correction forces escape to infinity. Applying
these results gives both directions of the new equivalence between bounded
correction and unboundedness.

Finally, $c_t(N)$ tends to its finite supremum and $1/x_t$ tends to zero.
The identity
\[
\frac{2^t}{3^{j_t}}=\bigl(N+c_t(N)\bigr)\frac1{x_t}
\]
proves the last limit in Theorem 1. Equivalently, on the additive logarithmic
scale, $j_t\log3-t\log2$ tends to positive infinity. The Lean declaration
uses the inverse-drift limit directly. Dividing by $t$ need not leave a
positive limit: no uniform strict density gap is claimed.

\section{Quantitative tails and the remaining absolute-error gap}
Define the closed-form tail budget
\[
B_K=\frac{1728}{5}\left(\frac{27}{32}\right)^K
 +\frac{7776}{13}\left(\frac{243}{256}\right)^K.
\]
It satisfies $B_K=b_K+B_{K+1}$, is nonnegative, and tends to zero.
\begin{theorem}
For any finite set $S$ of values on an unbounded orbit with $n\ge32^K$
for every $n\in S$, one has $\sum_{n\in S}1/n\le B_K$.
If a positive unbounded orbit stays above $32^K$ at all times, then
\[
\sum_{t\ge0}\frac1{x_t}\le B_K,
\qquad c_\infty(N)\le N\bigl(e^{B_K}-1\bigr).
\]
\end{theorem}
For a finite set below $32^{K+d}$, induction on $d$ proves the sharper bound
$\sum_{n\in S}1/n\le B_K-B_{K+d}$. Split the last shell and use
$B_{K+d}=b_{K+d}+B_{K+d+1}$. Every finite set fits below some such endpoint;
dropping the nonnegative remainder proves the first claim. Injectivity in
time transfers it to reciprocal partial sums. The earlier exponential
correction estimate
\[
N+c_t(N)\le N\exp\!\left(\sum_{s<t}\frac1{x_s}\right)
\]
then bounds every finite correction by $N(e^{B_K}-1)$; take the supremum.

There is also an exact relative tail identity, proved in
\texttt{RelativeCorrection.lean}:
\[
\frac{c_\infty(x_t)}{x_t}
 =\frac{c_\infty(N)-c_t(N)}{N+c_t(N)}\longrightarrow0.
\]
Its denominator has a positive limit. This statement gives relative
approximation, not a vanishing absolute error.

The explicit upper envelope cannot repair that distinction. Lean also proves
\[
32^K\bigl(e^{B_K}-1\bigr)\ge\frac{1728}{5}\,27^K.
\]
Indeed $e^u-1\ge u$ and the first summand of $B_K$ already gives the
right-hand side. Thus even at the smallest allowed seed, this upper bound
grows with $K$. This is a lower bound on the \emph{upper envelope}, not on
the actual correction. Independently, the existing ideal-barrier theorem
places the actual shifted correction strictly above one. Neither relative
convergence nor the quantitative envelope identifies an orbit value by
rounding the canonical ideal.

\section{Formal verification and provenance}
{\raggedright
The modules are \texttt{Collatz/OrbitPacking.lean},
\texttt{Collatz/OrbitSummability.lean}, and
\texttt{Collatz/RelativeCorrection.lean}. Principal declarations include
\texttt{unbounded\_orbit\_reciprocal\_summable},
\texttt{unbounded\_orbit\_supercritical},
\texttt{supercritical\_iff\_unbounded\_orbit}, and
\texttt{unbounded\_inverse\_drift\_tendsto\_zero}.
The quantitative declarations include
\texttt{unbounded\_reciprocal\_value\_tail},
\texttt{unbounded\_reciprocal\_tsum\_le},
\texttt{unbounded\_idealLimit\_le\_tailBudget}, and
\texttt{tail\_correction\_envelope\_lower}.
The build and selected audit pass: 120 declarations use only standard
foundations. Three exact finite controls check full small residue images,
the geometric identity, and the coalescence/multiplicity boundary.
The audit is dependency inspection, not independent kernel replay.\par}

A literature check prompted this route. Garcia and Tal,
\emph{A note on the generalized $3n+1$ problem}, Acta Arithmetica 90(3)
(1999), 245--250,
\href{https://doi.org/10.4064/aa-90-3-245-250}{doi:10.4064/aa-90-3-245-250},
give a power-saving interval estimate in equation (6), with orbit inclusion
in Corollary 1. That quantitative estimate implies reciprocal summability;
Banach density zero alone would not suffice. A July 2026
\href{https://mathoverflow.net/questions/513539}{MathOverflow question}
points out this implication. We checked the original six-page paper.
The present Lean proof uses a simpler origin-based packing argument with
explicit constants from this repository, importing neither Garcia--Tal's
theorem nor its cited Heppner input.

\section{What remains}
The earlier ledger's possible unbounded orbits with divergent reciprocal
series are now excluded. All hypothetical unbounded positive orbits lie
within the bounded-correction framework. This gives no contradiction within
that framework, no seed-independent absolute bound on all shifted ideal
corrections, and no exclusion of nontrivial cycles. The full Collatz
conjecture remains unresolved.
\end{document}
